\documentclass[../main]{subfiles}
\begin{document}
\section{Práctica de laboratorio motor paso a paso}
\begin{center}
  \img{images/10/steppers.jpg}{Steppers de impresora matriz de puntos EPSON}
\end{center}
\info{¿Qué es un motor paso a paso?}{Un motor paso a paso es un tipo de motor
  eléctrico que convierte señales eléctricas en movimientos precisos y
  controlados en forma de pasos. A diferencia de los motores convencionales
  que giran continuamente, los motores paso a paso se mueven en incrementos
  discretos o pasos.

  Los motores paso a paso se utilizan comúnmente en aplicaciones que requieren
  un control preciso de la posición, la velocidad y la aceleración, como en
  máquinas CNC, impresoras 3D, robots y equipos de automatización industrial.
  También se pueden encontrar en sistemas de seguimiento solar, control de
  satélites y sistemas de enfoque de cámaras.

  El funcionamiento de un motor paso a paso se basa en la activación secuencial
  de diferentes bobinas en el motor, que producen un campo magnético que hace
  girar el rotor en incrementos discretos. La velocidad y la dirección de
  rotación del motor se controlan mediante la secuencia de activación de las
  bobinas.
}
\actividad{Actividad práctica:}
\begin{enumerate}
  \item Elaborar un esquema o forma de identificar los terminales de cada
    motor paso a paso, con numeración o código de colores.
  \item Medir la resistencia entre todos los terminales de cada motor
    paso a paso, para esto realizar una tabla de doble entrada.
    \begin{center}
      \begin{tabular}{|c||c|c|c|c|c|c|c|c |}
        \hline
        \multicolumn{9}{|c|}{Tabla doble entrada} \\
        \hline\hline
        pin & 1 & 2 & 3 & 4 & 5 & 6 & 7 &8\\
        \hline\hline
        1 & & & & & & & &\\
        \hline
        2 & & & & & & & &\\
        \hline
        3 & & & & & & & &\\
        \hline
        4 & & & & & & & &\\
        \hline
        5 & & & & & & & &\\
        \hline
        6 & & & & & & & &\\
        \hline
        7 & & & & & & & &\\
        \hline
        8 & & & & & & & &\\
        \hline
      \end{tabular}
    \end{center}
  \item A partir de la tabla anterior, identificar bobinas en cada
    motor, y justificar.
  \item Prueba de freno: Conectar una bobina del motor a tensión de
    la fuente 1V, e intentar realizar el giro manual del eje del motor,
  \item Prueba led: Conectar un led a los terminales de la bobina del
    motor, y realizar el giro manual del eje.
  \item Clasificar cada motor según lo visto en la teoría, unipolar,
    bipolar, reluctancia variable, imán permanente,etc.
  \item Contar la cantidad de pasos, que posee el motor y calcular el
    ángulo de precisión de cada paso.
  \item Elaborar un \bf{informe} de las actividades realizadas con
    fotos y enviar a la plataforma classroom.
\end{enumerate}
\end{document}
